\begin{table}[p]
\centering
\caption{Selected historical sanity cases}
\label{tab:appendix_historical_sanity}
\small
\setlength{\tabcolsep}{4pt}
\renewcommand{\arraystretch}{1.08}
\begin{tabularx}{\textwidth}{@{}>{\raggedright\arraybackslash}p{0.24\textwidth}>{\raggedright\arraybackslash}p{0.22\textwidth}>{\raggedright\arraybackslash}p{0.14\textwidth}>{\raggedleft\arraybackslash}p{0.095\textwidth}>{\raggedleft\arraybackslash}p{0.105\textwidth}Y@{}}
\toprule
Case & City label & Genre family & First band & Emerg. & Geography note \\
\midrule
Birmingham heavy metal & Birmingham, United Kingdom & heavy metal & 1970 & 1977 & City-like label kept in baseline. \\
Tampa death metal & Tampa, United States & death metal & 1984 & 1987 & City-like label kept in baseline. \\
Bergen black metal & Bergen, Norway & black metal & 1991 & 1992 & City-like label kept in baseline. \\
Oslo black metal & Oslo, Norway & black metal & 1987 & 1990 & City-like label kept in baseline. \\
Gothenburg melodic death metal & Gothenburg, Sweden & melodic death metal & 1990 & 1996 & City-like label kept in baseline. \\
Bay Area thrash metal & Bay Area, United States & thrash metal & 1985 & 1992 & Region-like label excluded from city baseline. \\
\bottomrule
\end{tabularx}

\vspace{0.4em}
\parbox{0.96\textwidth}{\footnotesize Notes: These are operational scene dates from the project's
\texttt{city\_genre\_first\_appearance.csv} file, not claims about the exact historical birth date
of a genre. The point is narrower: the dates should be broadly plausible for canonical scenes, and
the geography note should clarify when a famous scene label is intentionally excluded because it is
a wider unit rather than a city-like observation.}
\end{table}
